\subsection{Additional details for high-dimensional Cox experiments}
\label{app:cox-device-exact}

To assess whether Shapley attribution remains practical when the number of
features greatly exceeds the number of patients, we consider two genomic
survival datasets: GSE24080 multiple-myeloma gene expression
\cite{geo_gse24080} and TCGA lower-grade glioma (LGG) DNA methylation
\cite{tcga_lgg_methylation}.
Genome-wide assays measure many molecular variables per patient, whereas
the cohorts with matched survival follow-up contain only hundreds of patients.
After eligibility filtering, GSE24080 contains 553 patients and 54,675
expression probe sets, and TCGA LGG contains 511 patients and 396,065
methylation sites. These are probe- or site-level features, retained without
gene aggregation. We use the supplied 339/214 training/validation split
for GSE24080 and a 383/128 split stratified by event status for TCGA LGG
(seed 42), giving an extreme small-$n$, large-$d$ setting.

We model overall survival using a Cox proportional-hazards model, which
incorporates right-censored observations: patients whose death has not been
observed by the end of follow-up. Its relative hazard factors across features,
making QuadraSHAP applicable directly:
\[
h(t\mid x)=h_0(t)f(x),
\qquad
f(x)=\exp(\beta^\top x)=\prod_{j=1}^{d}\exp(\beta_jx_j).
\]
To stabilize estimation when $d\gg n$, we shrink the Cox coefficients
towards zero using a ridge penalty, with the implementation's penalty
parameter fixed a priori at $\alpha=0.01d$. We use the Breslow approximation
for patients with equal recorded event times. Missing values are imputed using
training medians, and centering and scaling use training statistics.
All retained features have nonzero fitted coefficients.

We explain the relative hazard using 30 training-background patients,
sampled without replacement with seed 42 and held fixed across methods.
Specifically, absent features are replaced by the corresponding entries
of each background row and predictions are averaged:
\[
v_x(S)=\frac{1}{30}\sum_{b=1}^{30}
f\bigl(x_S,z^{(b)}_{\overline S}\bigr).
\]
The resulting game is a weighted sum of product games. We explain the
first five held-out patients in each saved split with the degree-exact rule
$m_q=\lceil d/2\rceil$ and with patient-specific node counts selected
to certify absolute quadrature tolerances
$\varepsilon\in\{10^{-1},10^{-2},10^{-3}\}$.

\paragraph{Matched degree-exact references.}
To assess approximation accuracy without changing the numerical backend,
we use a separate degree-exact reference for each device/precision setting.
Write $\widehat\phi^{a}_{m_p,p}$ for patient $p$'s computed attribution
vector using $m_p$ nodes, with $a=\mathrm{GPU}$ denoting GPU FP32 and
$a=\mathrm{CPU}$ denoting CPU FP64. Let $m_\star=\lceil d/2\rceil$.
The table reports
\[
\Delta_a=\max_{p=1,\ldots,5}
\|\widehat\phi^{a}_{m_p,p}-\widehat\phi^{a}_{m_\star,p}\|_\infty.
\]
Thus GPU approximation is compared only with GPU FP32 degree-exact output,
and CPU approximation only with CPU FP64 degree-exact output. These are
full degree-exact rules with 27,338 and 198,033 nodes, respectively; the
CPU references are newly measured rather than replaced with smaller-node
FP64 proxies. A patient passes if its maximum absolute feature-wise
discrepancy is at most its requested tolerance. Exact-row discrepancies
are zero by self-comparison and do not assert zero rounding error.

For the exact real-arithmetic Shapley vector $\phi_p$ and its quadrature
approximation $\phi_{m_p,p}$, the analytical node selector instead guarantees
\[
\|\phi_{m_p,p}-\phi_p\|_\infty\le C_p\le\varepsilon,
\qquad C_{\max}=\max_{p=1,\ldots,5}C_p.
\]
This bound can be conservative: the actual quadrature error need not be
close to either $C_p$ or $\varepsilon$. However, both computed vectors in
$\Delta_a$ also contain floating-point error. Consequently, a discrepancy
above $\varepsilon$ need not contradict the quadrature certificate, while
a discrepancy below $\varepsilon$ does not independently establish an
absolute error guarantee for either computed vector.

To avoid choosing unnecessarily many nodes, the implementation searches
integer node counts $m=1,2,\ldots$ and returns the first whose analytical
bound meets the requested tolerance; the degree-exact threshold is the
fallback cap. Evaluating the bound for each candidate requires optimizing
an auxiliary scalar that controls the bound. This inner optimization solves
its scalar stationarity equation by bisection, repeatedly halving an interval
containing the root. The separate \texttt{closed\_form\_budget} calculation
is recorded as a diagnostic and does not select the node count used in
these experiments. Timed automatic node selection includes both the
patient/background factor-summary pass and this integer search.

\paragraph{Implementation and timing.}
GPU explanations use JAX-Metal FP32; CPU explanations use JAX with FP64
enabled. Both evaluate the same parallel product-game formula, use an
8\,GiB memory-planning budget and one background row per block, and
accumulate returned blocks on the host in FP64. The memory setting is a
planning budget, not a measurement of actual allocation.
Each table entry summarizes five distinct held-out patients, with one
timed explanation per patient and method. We report the arithmetic mean
and sample standard deviation across those patients. This variation is
not repeated-run uncertainty and does not establish statistical significance.

GPU exact and approximate measurements are reused from the matched saved
experiments, and CPU approximate timings reuse the warmed CPU FP64 sweep.
The full degree-exact CPU runs are new. Model coefficients, preprocessing,
background and held-out rows, numerical source hashes, and package versions
are verified to match. Timings exclude rule construction and recorded
warm-up. CPU degree-exact computation uses a checkpointed integration loop
with the same factors, node order, and background/block accumulation order
as the public API; a small synthetic comparison was bitwise identical.
Before each dataset, one actual node block with the full exact-core shape
is evaluated for warm-up, followed by three pilot core calls. These calls
are excluded; the warm-up is not a full degree-exact explanation.
The reported CPU exact time sums active integration, including factor
preparation, block planning, transfers, synchronized core evaluation, and
host accumulation. Checkpoint I/O is measured separately and excluded;
if resumed, inactive downtime and uncheckpointed work lost to interruption
are also excluded. Thus this timing is not an undifferentiated process
wall-clock measurement.
Approximate API timings include automatic node selection as well as
explanation, synchronization, transfers and host work. Measurements were
recorded in separate runs. GPU FP32 versus CPU FP64 changes both device
and precision, so their latency ratio is not a matched-precision GPU speedup.

The quadrature nodes and weights are cached across patients. Previously
measured CPU construction of the exact GSE24080 and TCGA rules required
10.97\,s and 561.92\,s, respectively; the latest approximate rules required
0.386--0.563\,ms (range of medians over five fresh constructions per rule).
These setup costs are separate from explanation latency. Caching rules
does not remove evaluations of the patient- and background-dependent
integrands at those nodes.

\paragraph{Numerical interpretation and scope.}
All 30 selected approximate rules satisfy $C_p\le\varepsilon$.
GPU FP32 meets its same-backend discrepancy threshold in 26/30
cases, while CPU FP64 meets it in 30/30 cases.
The largest CPU discrepancy across patients, features, and the three
tolerances is $2.44\times10^{-7}$.
The CPU FP64 results provide empirical evidence that the selected rules
agree closely with full degree-exact quadrature at higher precision.
They must be paired with the CPU runtime, not assigned to GPU FP32 timing.
The numerical-reference convention in this table is distinct from earlier
checks against a shared, independently implemented 64-node FP64 reference.
Those checks and the separate CPU FP32 control are retained as diagnostic
artifacts; they are not the definitions of either discrepancy reported here.

Finally, the explanations concern the fitted relative-hazard model and the
fixed uniform distribution on the 30 selected background rows. The
certificate does not quantify model-estimation uncertainty, uncertainty
from choosing a finite background sample, or discrepancy from a population
background distribution. The survival datasets demonstrate computational
feasibility in a small-$n$, large-$d$ setting; these experiments do not
establish clinical usefulness or causal interpretation of the attributions.
